\chapter{Vettori tangenti a varietà}
    \section{Vettori tangenti a spazio tangente}
        Sia $M$ varietà differenziabile di dimensione $n$ e
        sia $p \in M$ consideriamo l'insieme
        \begin{align*}
        \mathcal{F}(p) &= \left\{ f : A \to \mathbb{R} \;\middle|\; A \text{ è intorno aperto di } M \text{ in } p \text{ e } f \text{ è differenziabile} \right\} \\
        &= \bigcup_{\substack{A \text{ intorno} \\ \text{aperto di } p}} \mathcal{F}(A).
        \end{align*}

        Su $\mathcal{F}(p)$ si possono definire tre leggi di composizione:
        \begin{enumerate}
        \item $+ : \mathcal{F}(p) \times \mathcal{F}(p) \to \mathcal{F}(p)$
        \item $\cdot : \mathcal{F}(p) \times \mathcal{F}(p) \to \mathcal{F}(p)$
        \item $\cdot : \mathbb{R} \times \mathcal{F}(p) \to \mathcal{F}(p)$
        \end{enumerate}
        tali che $\forall \lambda \in \mathbb{R}$, $\forall f, g \in \mathcal{F}(p)$, con $f : A \to \mathbb{R}$ e $g : B \to \mathbb{R}$ funzioni differenziabili su $A$ e $B$ intorni aperti di $p$:
        \begin{align*}
        f + g &: A \cap B \to \mathbb{R}, & q \in A \cap B &\longmapsto f(q) + g(q) \\
        f \cdot g &: A \cap B \to \mathbb{R}, & q \in A \cap B &\longmapsto f(q) \cdot g(q) \\
        \lambda \cdot f &: A \to \mathbb{R}, & q \in A &\longmapsto \lambda \cdot f(q)
        \end{align*}

        Le due leggi di composizione interna $+$ e $\cdot$ sono associative e commutative.
        L'addizione ammette elemento neutro dato da
        \[
        f_0 : M \to \mathbb{R}, \qquad q \longmapsto f_0(q) = 0.
        \]
        $\forall f \in \mathcal{F}(p)$, con $f : A \to \mathbb{R}$ differenziabile su $A$ intorno aperto di $p$, risulta $f_0 + f = f = f + f_0$.

        Non tutti gli elementi di $\mathcal{F}(p)$ ammettono opposto. Infatti, se $f : A \to \mathbb{R}$ con $A$ intorno aperto di $p$, considerata $-f : A \to \mathbb{R}$, $q \longmapsto -f(q)$, risulta $f + (-f) = f_0$ solo se $A = M$.

        \begin{definizione}
        \mbox{}

        Sia $M$ varietà differenziabile di dimensione $n$. Sia $\textcolor{red}{p \in M}$. Si definisce \textcolor{red}{vettore tangente} in $p$ a $M$ ogni applicazione $\textcolor{red}{v : \mathcal{F}(p) \to \mathbb{R}}$ tale che
        \begin{enumerate}
        \item $\forall f, g \in \mathcal{F}(p) \quad \textcolor{red}{v(f + g) = v(f) + v(g)}$
        \item $\forall \lambda \in \mathbb{R} \; \forall f \in \mathcal{F}(p) \quad \textcolor{red}{v(\lambda f) = \lambda v(f)}$
        \item $\forall f, g \in \mathcal{F}(p) \quad \textcolor{red}{v(f \cdot g) = f(p)v(g) + v(f) \cdot g(p)}$
        \end{enumerate}

        L'insieme dei vettori tangenti in $p$ a $M$ si denota con $\textcolor{red}{T_p M}$ e prende il nome di \textcolor{red}{spazio tangente} in $p$ a $M$:
        \[
        \textcolor{red}{T_p M = \{ v : \mathcal{F}(p) \to \mathbb{R} \mid v \text{ è vettore tangente in } p \text{ a } M \}}.
        \]
        \end{definizione}

        \begin{proposizione}
        \mbox{}

        $\textcolor{red}{T_p M}$ è uno spazio vettoriale reale.
        \end{proposizione}

        \begin{proof}
        L'insieme $V = \{ v : \mathcal{F}(p) \to \mathbb{R} \}$ è uno spazio vettoriale reale, munito delle leggi di composizione
        \[
        + : V \times V \to V, \qquad \cdot : \mathbb{R} \times V \to V
        \]
        tali che $\forall u, v \in V$, $\forall \lambda \in \mathbb{R}$:
        \begin{align*}
        u + v &: \mathcal{F}(p) \to \mathbb{R}, & f &\longmapsto u(f) + v(f) \\
        \lambda \cdot v &: \mathcal{F}(p) \to \mathbb{R}, & f &\longmapsto \lambda \cdot v(f)
        \end{align*}
        Si verifica che $T_p M$ è sottospazio vettoriale di $V$, ovvero:
        \[
        \left.
        \begin{aligned}
        &\forall u, v \in T_p M \quad u + v \in T_p M \\
        &\forall \lambda \in \mathbb{R} \; \forall v \in T_p M \quad \lambda v \in T_p M
        \end{aligned}
        \right\} \text{ esercizio.}
        \]
        \end{proof}

        \begin{proposizione}
        \mbox{}

        Sia $M$ varietà differenziabile di dimensione $n$. Siano $p \in M$, $\textcolor{red}{u, v \in T_p M}$.
        \begin{enumerate}
        \item Se $U$ è intorno aperto di $p$ e $c : U \to \mathbb{R}$ è funzione costante, allora $v(c) = 0$.
        \item Se $f, g \in \mathcal{F}(p)$ coincidono in un intorno aperto di $p$, allora $v(f) = v(g)$.
        \item Se $\forall f \in \mathcal{F}(M) \; u(f) = v(f)$, allora $u = v$.
        \item $v = 0 \iff \forall f \in \mathcal{F}(M) \; v(f) = 0$ \textcolor{blu}{(segue immediatamente dalla 3)}.
        \end{enumerate}
        \end{proposizione}

        \begin{proof}
        \begin{enumerate}
        \item Sia $c : U \to \mathbb{R}$ tale che $\forall q \in U \; c(q) = c \in \mathbb{R}$, e sia $c_1 : U \to \mathbb{R}$ tale che $\forall q \in U \; c_1(q) = 1$. Allora:
        \begin{align*}
        v(c \cdot c_1) &= c(p) \cdot v(c_1) + v(c) \cdot c_1(p) \\
        \implies c \cdot v(c_1) &= c \cdot v(c_1) + v(c) \implies v(c) = 0.
        \end{align*}

        \item Siano $f, g \in \mathcal{F}(p)$, con $f : A \to \mathbb{R}$ e $g : B \to \mathbb{R}$, $A, B$ intorni aperti di $p$. Sia $U$ intorno aperto di $p$ tale che $U \subset A \cap B$ e $f|_U = g|_U$. Considerata $c_1 : U \to \mathbb{R}$, $q \longmapsto 1$, si ha:
        \[
        \left.
        \begin{aligned}
        c_1 \cdot f &: \overbrace{U}^{= U \cap A} \to \mathbb{R}, & c_1 \cdot f &= f|_U \\
        c_1 \cdot g &: \overbrace{U}^{= U \cap B} \to \mathbb{R}, & c_1 \cdot g &= g|_U
        \end{aligned}
        \right\} \implies c_1 \cdot f = c_1 \cdot g.
        \]
        Calcolando $v$:
        \begin{align*}
        v(c_1 \cdot f) &= c_1(p) \cdot v(f) + \underbrace{v(c_1)}_{= 0} \cdot f(p) = v(f) \\
        v(c_1 \cdot g) &= c_1(p) \cdot v(g) + v(c_1) \cdot g(p) = v(g)
        \end{align*}
        Pertanto $v(f) = v(g)$.

        \item Siano $u, v \in T_p M$ tali che $\textcolor{orange}{\forall f \in \mathcal{F}(M) \; u(f) = v(f)}$. Proviamo che $\forall f \in \mathcal{F}(p) \; u(f) = v(f)$.

        Sia $f \in \mathcal{F}(p)$. Pertanto $f : A \to \mathbb{R}$ differenziabile, $A$ intorno aperto di $p$. Esiste $W$ intorno aperto di $p$, $W \subset A$, ed esiste $\tilde{f} : M \to \mathbb{R}$ differenziabile tale che $\tilde{f}|_W = f|_W$ e $\supp(\tilde{f}) \subset A$.

        Poiché $\tilde{f}, f \in \mathcal{F}(p)$ coincidono su un intorno aperto di $p$, per il punto 2 si ha:
        \[
        u(\tilde{f}) = u(f) \quad \text{e} \quad v(\tilde{f}) = v(f).
        \]
        Per ipotesi $\textcolor{orange}{u(\tilde{f}) = v(\tilde{f})}$. Pertanto $u(f) = v(f)$.
        \end{enumerate}
        \end{proof}


    \section{Vettori tangenti coordinati}

        Sia $M$ varietà differenziabile di dimensione $n$. Sia $p \in M$. Sia $(U, \varphi)$ carta locale tale che $p \in U$, con $\varphi : U \to \varphi(U) \subset \mathbb{R}^n$ diffeomorfismo.
        Siano $x^1, \dots, x^n : U \to \mathbb{R}$ le funzioni coordinate locali della carta $(U, \varphi)$: $x^i = p_i \circ \varphi : U \to \mathbb{R}$.

        \begin{definizione}
        \mbox{}

        Sia $i = 1, \dots, n$ e definiamo un'applicazione
        \[
        \textcolor{red}{\left( \frac{\partial}{\partial x^i} \right)_p : \mathcal{F}(p) \to \mathbb{R}}
        \]
        nel modo seguente: se $f \in \mathcal{F}(p)$, allora $f : A \to \mathbb{R}$ è funzione differenziabile su $A$ intorno aperto di $p$. Per definizione di funzione differenziabile, considerata la carta locale $(U, \varphi)$, la funzione
        \[
        f \circ \varphi^{-1} : \varphi(U \cap A) \subset \mathbb{R}^n \to \mathbb{R} \qquad \left( \varphi(U \cap A) \xrightarrow{\varphi^{-1}} U \cap A \xrightarrow{f} \mathbb{R} \right)
        \]
        è differenziabile. Poniamo quindi:
        \[
        \textcolor{red}{\left( \frac{\partial}{\partial x^i} \right)_p (f) = \frac{\partial (f \circ \varphi^{-1})}{\partial u^i} (\varphi(p))}
        \]
        \textcolor{blu}{(derivata parziale di $f \circ \varphi^{-1}$ rispetto alla variabile $u^i$ nel punto $\varphi(p)$)}.
        \end{definizione}

        \begin{proposizione}
        \mbox{}

        $\forall i = 1, \dots, n \quad \textcolor{red}{\left( \frac{\partial}{\partial x^i} \right)_p} \in T_p M$.
        \end{proposizione}

        \begin{proof}
        Bisogna verificare che $\forall f, g \in \mathcal{F}(p)$ e $\forall \lambda \in \mathbb{R}$:
        \begin{enumerate}
        \item $\left( \frac{\partial}{\partial x^i} \right)_p (f + g) = \left( \frac{\partial}{\partial x^i} \right)_p (f) + \left( \frac{\partial}{\partial x^i} \right)_p (g)$
        \item $\left( \frac{\partial}{\partial x^i} \right)_p (\lambda f) = \lambda \left( \frac{\partial}{\partial x^i} \right)_p (f)$
        \item $\left( \frac{\partial}{\partial x^i} \right)_p (f \cdot g) = f(p) \left( \frac{\partial}{\partial x^i} \right)_p (g) + \left( \frac{\partial}{\partial x^i} \right)_p (f) \cdot g(p)$
        \end{enumerate}

        Verifichiamo la 3:
        \begin{align*}
        \left( \frac{\partial}{\partial x^i} \right)_p (f \cdot g) &= \frac{\partial ((f \cdot g) \circ \varphi^{-1})}{\partial u^i} (\varphi(p)) \\
        &= \frac{\partial ((f \circ \varphi^{-1}) \cdot (g \circ \varphi^{-1}))}{\partial u^i} (\varphi(p)) \\
        &= (f \circ \varphi^{-1})(\varphi(p)) \cdot \frac{\partial (g \circ \varphi^{-1})}{\partial u^i} (\varphi(p)) + \frac{\partial (f \circ \varphi^{-1})}{\partial u^i} (\varphi(p)) \cdot (g \circ \varphi^{-1})(\varphi(p)) \\
        &= f(p) \cdot \left( \frac{\partial}{\partial x^i} \right)_p (g) + \left( \frac{\partial}{\partial x^i} \right)_p (f) \cdot g(p).
        \end{align*}
        \end{proof}

        \begin{definizione}
        \mbox{}

        Se $(U, \varphi)$ è una carta locale di $M$ in $p$, di coordinate locali $x^1, \dots, x^n : U \to \mathbb{R}$, e se $f \in \mathcal{F}(p)$, porremo
        \[
        \textcolor{red}{\left( \frac{\partial f}{\partial x^i} \right)_p := \left( \frac{\partial}{\partial x^i} \right)_p (f)} \qquad i = 1, \dots, n
        \]
        detta \textcolor{red}{derivata parziale} di $f$ in $p$ rispetto alla $i$-esima coordinata locale $x^i$ della carta $(U, \varphi)$.
        \end{definizione}

        \begin{osservazione}
        \mbox{}

        Sia $(U, \varphi)$ una carta locale di $M$ di coordinate locali $x^1, \dots, x^n : U \to \mathbb{R}$. Sia $p \in U$. Per ogni $j = 1, \dots, n$, $x^j \in \mathcal{F}(p)$. Si può quindi applicare il vettore tangente $\left( \frac{\partial}{\partial x^i} \right)_p$ alla funzione $x^j$:
        \begin{align*}
        \left( \frac{\partial x^j}{\partial x^i} \right)_p &= \left( \frac{\partial}{\partial x^i} \right)_p (x^j) = \frac{\partial (x^j \circ \varphi^{-1})}{\partial u^i} (\varphi(p)) \\
        &= \frac{\partial (p_j \circ \varphi \circ \varphi^{-1})}{\partial u^i} (\varphi(p)) \\
        &= \frac{\partial p_j}{\partial u^i} (\varphi(p)) = \frac{\partial u^j}{\partial u^i} (\varphi(p)) = \delta_{ij}.
        \end{align*}
        \end{osservazione}

        \begin{lemma}[Dimensione dello spazio tangente]
        \mbox{}

        Sia $M$ varietà differenziabile, $\dim M = n \geq 1$. Siano $p \in M$, $A$ intorno aperto di $p$, $f : A \to \mathbb{R}$ funzione differenziabile. Sia $(U, \varphi)$ carta locale di $M$, $p \in U$. Siano $x^1, \dots, x^n : U \to \mathbb{R}$ le funzioni coordinate locali. Poniamo $\varphi(p) = (x^1(p), \dots, x^n(p)) = (c^1, \dots, c^n) = c$.
        Allora esiste $W$ intorno aperto di $p$ tale che $W \subset U \cap A$, ed esistono $n$ funzioni differenziabili $f_1, \dots, f_n : W \to \mathbb{R}$ tali che
        \[
        \forall q \in W \quad f(q) = f(p) + \sum_{i=1}^n (x^i(q) - c^i) f_i(q) \qquad \text{e} \qquad f_i(p) = \left( \frac{\partial f}{\partial x^i} \right)_p.
        \]
        \end{lemma}

        \begin{teorema}[Dimensione dello spazio tangente]
        \mbox{}

        Sia $M$ varietà differenziabile di dimensione $n \geq 1$. Sia $p \in M$. Allora $\dim T_p M = n$.
        In particolare, se $(U, \varphi)$ è una carta locale, $p \in U$, allora
        \[
        \textcolor{red}{\left\{ \left( \frac{\partial}{\partial x^1} \right)_p, \dots, \left( \frac{\partial}{\partial x^n} \right)_p \right\}}
        \]
        è \textcolor{red}{base} di $T_p M$, detta \textcolor{red}{base coordinata} di $T_p M$ associata alla carta $(U, \varphi)$.
        \end{teorema}

        \begin{proof}
        Sia $(U, \varphi)$ carta locale, $p \in U$. Siano $x^1, \dots, x^n : U \to \mathbb{R}$ funzioni coordinate locali. Consideriamo i vettori tangenti a $M$ in $p$
        \[
        \left( \frac{\partial}{\partial x^1} \right)_p, \dots, \left( \frac{\partial}{\partial x^n} \right)_p
        \]
        e proviamo che formano una base di $T_p M$.

        \begin{itemize}
        \item \textbf{Lineare indipendenza:} Siano $a^1, \dots, a^n \in \mathbb{R}$ tali che
        \[
        \underbrace{\sum_{i=1}^n a^i \left( \frac{\partial}{\partial x^i} \right)_p}_{=: \, v} = 0.
        \]
        Allora $\forall f \in \mathcal{F}(p) \; v(f) = 0$. In particolare, $\forall j = 1, \dots, n$, $x^j \in \mathcal{F}(p) \implies v(x^j) = 0$, cioè:
        \[
        \sum_{i=1}^n a^i \left( \frac{\partial}{\partial x^i} \right)_p (x^j) = 0 \implies \sum_{i=1}^n a^i \delta_{ij} = 0 \implies a^j = 0.
        \]

        \item \textbf{Proviamo che i vettori generano $T_p M$:} Sia $v \in T_p M$, $v : \mathcal{F}(p) \to \mathbb{R}$. Sia $f \in \mathcal{F}(p)$, con $f : A \to \mathbb{R}$ differenziabile, $A$ intorno aperto di $p$.
        Per il lemma, $\exists W$ intorno aperto di $p$, $W \subset U \cap A$, ed $\exists f_1, \dots, f_n : W \to \mathbb{R}$ funzioni differenziabili tali che
        \[
        \forall q \in W \quad f(q) = f(p) + \sum_{i=1}^n f_i(q) (x^i(q) - c^i) \qquad \text{e} \qquad f_i(p) = \left( \frac{\partial f}{\partial x^i} \right)_p.
        \]
        Pertanto, $f|_W = f(p) + \sum_{i=1}^n f_i (x^i|_W - c^i)$. Applicando $v$ (e ricordando che $v$ coincide sul germe di $f|_W$ e che $v$ si annulla sulle costanti):
        \begin{align*}
        v(f) &= \sum_{i=1}^n \left[ f_i(p) \, v(x^i - c^i) + v(f_i) (x^i(p) - c^i) \right] \\
        &= \sum_{i=1}^n \left( \frac{\partial f}{\partial x^i} \right)_p v(x^i) = \sum_{i=1}^n v(x^i) \left( \frac{\partial}{\partial x^i} \right)_p (f).
        \end{align*}
        Pertanto:
        \[
        v = \sum_{i=1}^n v(x^i) \left( \frac{\partial}{\partial x^i} \right)_p.
        \]
        \end{itemize}
        Si noti che la $i$-esima componente di $v$ nella base coordinata è data da $v(x^i)$. In particolare:
        \[
        v = 0 \iff \forall i = 1, \dots, n \quad v(x^i) = 0.
        \]
        \end{proof}


    \section{Matrici di passaggio tra basi coordinate}

        Sia $M$ varietà differenziabile di dimensione $n \geq 1$. Sia $p \in M$.
        Siano $(U, \varphi)$, $(U', \varphi')$ carte locali di $M$ tali che $p \in U \cap U'$.
        Siano $x^1, \dots, x^n : U \to \mathbb{R}$ coordinate locali di $(U, \varphi)$ e $x'^1, \dots, x'^n : U' \to \mathbb{R}$ coordinate locali di $(U', \varphi')$, con
        \[
        x^i = p_i \circ \varphi, \qquad x'^i = p_i \circ \varphi' \qquad \forall i = 1, \dots, n.
        \]
        Siano
        \[
        \mathcal{B} = \left\{ \left( \frac{\partial}{\partial x^i} \right)_p \right\}_{1 \leq i \leq n}, \qquad \mathcal{B}' = \left\{ \left( \frac{\partial}{\partial x'^i} \right)_p \right\}_{1 \leq i \leq n}
        \]
        le basi coordinate di $T_p M$ relative alle carte $(U, \varphi)$ e $(U', \varphi')$.

        Determiniamo la matrice di passaggio dalla base $\mathcal{B}'$ alla base $\mathcal{B}$:
        \[
        \forall j = 1, \dots, n \qquad \left( \frac{\partial}{\partial x^j} \right)_p = \sum_{i=1}^n a^i_j \left( \frac{\partial}{\partial x'^i} \right)_p.
        \]
        Per quanto visto precedemente, la $i$-esima componente rispetto alla base $\mathcal{B}'$ è data dalla valutazione su $x'^i$:
        \begin{align*}
        a^i_j &= \left( \frac{\partial}{\partial x^j} \right)_p (x'^i) = \left( \frac{\partial x'^i}{\partial x^j} \right)_p \\
        &= \frac{\partial (x'^i \circ \varphi^{-1})}{\partial u^j} (\varphi(p)) \\
        &= \frac{\partial (p_i \circ \varphi' \circ \varphi^{-1})}{\partial u^j} (\varphi(p)).
        \end{align*}
        Poiché $\varphi' \circ \varphi^{-1} : \varphi(U \cap U') \to \varphi'(U \cap U')$ è diffeomorfismo, $p_i \circ \varphi' \circ \varphi^{-1} : \varphi(U \cap U') \to \mathbb{R}$ è la $i$-esima funzione componente di $\varphi' \circ \varphi^{-1}$.

        Pertanto $a^i_j = \left( \frac{\partial x'^i}{\partial x^j} \right)_p$ è il coefficiente di posto $(i, j)$ della matrice Jacobiana della funzione di transizione $\varphi' \circ \varphi^{-1}$ nel punto $\varphi(p)$.

        Pertanto:
        \[
        \left( \frac{\partial}{\partial x^j} \right)_p = \sum_{i=1}^n \left( \frac{\partial x'^i}{\partial x^j} \right)_p \left( \frac{\partial}{\partial x'^i} \right)_p, \qquad M_{\mathcal{B}' \mathcal{B}} = J(\varphi' \circ \varphi^{-1})_{\varphi(p)}.
        \]


    \section{Spazio tangente di uno spazio vettoriale reale}

        \begin{proposizione}
        \mbox{}

        Sia $\textcolor{red}{V}$ spazio vettoriale reale, $\dim V = n \geq 1$, munito di struttura di varietà differenziabile di dimensione $n$. Allora $\forall b \in V$ esiste un isomorfismo canonico
        \[
        \Phi : T_b V \longrightarrow V.
        \]
        \end{proposizione}

        \begin{proof}
        \mbox{}
        
        Sia $\mathcal{B} = \{e_1, \dots, e_n\}$ base di $V$, e sia $\varphi : V \to \mathbb{R}^n$ l'isomorfismo tale che
        \[
        v = x^i e_i \longmapsto (x^1, \dots, x^n).
        \]
        Consideriamo su $V$ la struttura differenziabile definita dall'atlante $\{(V, \varphi)\}$ \textcolor{blu}{(è stato provato che la struttura non dipende dalla base $\mathcal{B}$)}.

        Le funzioni coordinate locali nella carta $(V, \varphi)$ sono $x^i = p_i \circ \varphi : V \to \mathbb{R}$, $v = x^k e_k \longmapsto x^i$.
        Sia $\left\{ \left( \frac{\partial}{\partial x^i} \right)_b \right\}_{1 \leq i \leq n}$ la base coordinata di $T_b V$ relativa alla carta $(V, \varphi)$.

        Sia $\Phi_{\mathcal{B}} : T_b V \to V$ l'isomorfismo tale che
        \[
        \forall i = 1, \dots, n \qquad \Phi_{\mathcal{B}} \left( \left( \frac{\partial}{\partial x^i} \right)_b \right) = e_i.
        \]
        Verifichiamo che l'isomorfismo $\Phi_{\mathcal{B}}$ non dipende dalla base $\mathcal{B}$, ed è quindi canonico.

        Sia $\mathcal{B}' = \{e'_1, \dots, e'_n\}$ un'altra base di $V$, e $(V, \varphi')$ la carta su $V$ definita da $\mathcal{B}'$, ovvero $\varphi' : V \to \mathbb{R}^n$, $v = x'^i e'_i \longmapsto (x'^1, \dots, x'^n)$.
        Sia $\left\{ \left( \frac{\partial}{\partial x'^i} \right)_b \right\}_{1 \leq i \leq n}$ la base coordinata di $T_b V$, e $\Phi_{\mathcal{B}'} : T_b V \to V$ l'isomorfismo tale che
        \[
        \forall i = 1, \dots, n \qquad \Phi_{\mathcal{B}'} \left( \left( \frac{\partial}{\partial x'^i} \right)_b \right) = e'_i.
        \]
        Dobbiamo verificare che $\Phi_{\mathcal{B}} = \Phi_{\mathcal{B}'}$.

        Sappiamo che
        \[
        \left( \frac{\partial}{\partial x^j} \right)_b = \sum_{i=1}^n \left( \frac{\partial x'^i}{\partial x^j} \right)_b \left( \frac{\partial}{\partial x'^i} \right)_b,
        \]
        dove $\left( \frac{\partial x'^i}{\partial x^j} \right)_b$ è il coefficiente di posto $(i, j)$ della matrice Jacobiana della funzione di transizione $\varphi' \circ \varphi^{-1}$ in $\varphi(b)$.

        Si ha $\varphi' \circ \varphi^{-1} = L_A : X \in \mathbb{R}^n \longmapsto AX$, dove $A = M_{\mathcal{B}', \mathcal{B}}$ (cioè $e_j = a^i_j e'_i$):
        \[
        X = (x^1, \dots, x^n) \xrightarrow{\varphi^{-1}} v = x^i e_i \xrightarrow{\varphi'} X' = (x'^1, \dots, x'^n) \quad \text{t.c. } v = x'^i e'_i, \; x'^i = a^i_j x^j.
        \]
        La matrice Jacobiana di $L_A$ è $A$. Pertanto:
        \[
        \left( \frac{\partial}{\partial x^j} \right)_b = a^i_j \left( \frac{\partial}{\partial x'^i} \right)_b.
        \]
        Applicando $\Phi_{\mathcal{B}'}$:
        \begin{align*}
        \Phi_{\mathcal{B}'} \left( \left( \frac{\partial}{\partial x^j} \right)_b \right) &= a^i_j \, \Phi_{\mathcal{B}'} \left( \left( \frac{\partial}{\partial x'^i} \right)_b \right) \\
        &= a^i_j e'_i = e_j = \Phi_{\mathcal{B}} \left( \left( \frac{\partial}{\partial x^j} \right)_b \right) \qquad \forall j = 1, \dots, n.
        \end{align*}
        Pertanto $\Phi_{\mathcal{B}'} = \Phi_{\mathcal{B}}$.
        \end{proof}